\section{代数簇}\label{sec:19}\label{section:19}
\subsection{多重次调和函数}\label{sec:19a}
Stein 流形是一个复流形 $(X,I_X)$，它容许一个穷竭的，也就是固有且有下界的多重次调和函数 $\phi:X\to\R$。多重次调和性是说，令 $\theta_\phi=-d^c\phi=-d\phi\circ I_X$ 时，$\omega_\phi=d\theta_\phi=-dd^c\phi$ 是 Kähler 形式。注意，$dd^c\phi$ 总是 $(1,1)$ 型的；非平凡的条件是，它在复切直线上为负，复切直线指 $TX$ 中带自然定向的实 $I_X$-不变二维子空间。一个直接推论是，给定任意全纯映射 $u:S\to X$，函数 $\rho=\phi\circ u$ 如\eref{7.4}那样满足 $\Delta\rho\ge0$，因此服从最大值原理。

由于多重次调和性条件具有凸性，比较容易对这类函数作调整。下面需要其中两个初等实例；有关这一方向的更多内容，见\bcite{42}或\bcite{22}。

\SourcePage{300}
\StatementNumber{1}\begin{lemma}\label{st:19.1}
设 $\phi$ 为穷竭的多重次调和函数，$h:\R\to\R$ 处处满足 $h'>0$、$h''\ge0$。则 $h(\phi)$ 仍为穷竭的多重次调和函数。
\end{lemma}
\begin{proof}
有 $dd^ch(\phi)=h'(\phi)dd^c\phi+h''(\phi)d\phi\wedge d^c\phi$。在复切直线上，第一个加项为负，第二个非正。\znote{292}{本书的 $d^c$ 不带 $\tfrac12$：由 $\theta_\phi=-d^c\phi=-d\phi\circ I_X$ 得 $d^c\phi=d\phi\circ I_X$，而不少文献取 $d^c=\tfrac i2(\bar\partial-\partial)$，换算 $\omega_\phi=-dd^c\phi$ 时会差一个因子 $2$。第二个加项非正可直接验证：在复切直线上取基 $v,Iv$，则 $d\phi\wedge d^c\phi(v,Iv)=-(d\phi(v)^2+d\phi(Iv)^2)\le0$，而 $h''\ge0$ 又使它无法抵消第一项的负性；若 $h$ 为仿射函数则第二项整个消失，结论同样成立（引理 19.2 证明中的 $l$ 在 $t\ll0$ 与 $t\gg0$ 两段就分别取仿射形式）。[来源：A，原书 p291 对 $\theta_\phi$、$d^c$ 的定义与引理 19.1、19.2 的证明。]}
\end{proof}

\StatementNumber{2}\begin{lemma}\label{st:19.2}
设 $\phi,\psi$ 为 $(X,I_X)$ 上穷竭的多重次调和函数。假设 $\psi$ 比 $\phi$ 增长得快，意思是 $\psi-\phi$ 仍为穷竭函数。固定 $X$ 的某个紧子集。那么，存在另一个穷竭的多重次调和函数 $\widetilde\psi$，它在该紧子集上等于 $\phi$，在无穷远处等于 $\psi+\text{常数}$。
\end{lemma}
\begin{proof}
选择函数 $l$，使 $t\ll0$ 时 $l(t)=0$，$t\gg0$ 时 $l(t)=t$，且处处 $l''(t)\ge0$。对某个常数 $b$，取 $\xi=\phi+l(\psi-\phi-b)$。由与前面相同的理由，它是多重次调和的：
\[
\begin{aligned}
dd^c\xi={}&\bigl(1-l'(\psi-\phi-b)\bigr)dd^c\phi
+l'(\psi-\phi-b)dd^c\psi\\
&+l''(\psi-\phi-b)d(\psi-\phi)\wedge d^c(\psi-\phi).
\end{aligned}
\]
将常数 $b$ 取充分大，即可保证在选定的紧子集上 $l(\psi-\phi-b)=0$。
\end{proof}

若 $\phi$ 的临界点集紧，就称它为\emph{有限型}。一般来说，判断一个 Stein 流形是否容许这样的函数可能很困难；对于四维流形，微分结构会带来微妙的障碍\bcite{63}。不过，在许多大类情形中，这很容易做到。

\StatementNumber{3}\begin{example}\label{st:19.3}
(i) 设 $\overline X$ 为光滑复射影簇，带有丰沛线丛 $\mathcal L$。丰沛性意味着该线丛上有度量 $|\cdot|$，其曲率是 $\overline X$ 上的 Kähler 形式。设 $s$ 为 $\mathcal L$ 的全纯截面，使 $s^{-1}(0)$ 是正规交叉除子。那么
\[
\phi(x)=-\log|s(x)|
\]
是 $X=\overline X\setminus s^{-1}(0)$ 上有限型的穷竭严格多重次调和函数。相应 Kähler 形式 $\omega_\phi$ 是前述 Kähler 形式在 $X$ 上的限制。有限型性质可由直接的局部计算证明，见\bcite{133}，引理~8。

(ii) 设 $X\subset\CC^N$ 为光滑仿射代数簇，$\phi:X\to\R$ 为穷竭的多重次调和函数，而且是一个实多项式的限制。那么 $\phi$ 自动是有限型的，因为其临界点集本身是实代数簇，因此只有有限多个连通分支，每个分支对应于唯一的临界值。

(iii) 类似地，设 $q:\CC^N\to\CC$ 为复多项式，原点是它的孤立临界点，且 $q(0)=0$、$Dq(0)=0$。那么，对
\SourcePage{301}
所有 $0<|\epsilon|\ll\delta$，Milnor 纤维 $X=\{x\in\CC^N:q(x)=\epsilon,\ |x|<\delta\}$ 是 Stein 流形，$\phi(x)=(\delta^2-|x|^2)^{-1}$ 是其上的有限型多重次调和函数。注意，Milnor 纤维一般不双全纯同构于代数簇。
\end{example}

此后固定 $(X,I_X)$，并要求它带有二次复体积形式 $\eta_X^2$ 及有限型的穷竭多重次调和函数 $\phi$。取 $c$ 大于 $\phi$ 的全部临界值，考虑次水平集
\[
M=X^{\phi,c}=\phi^{-1}(({-\infty},c])
\]
及 $\theta_X,\omega_X$ 的限制。这是带接触型边界的正合辛流形，$I_X$ 的限制也是边界附近的接触型殆复结构。再使用 $\eta_X^2$ 的限制，就具备了定义 Fukaya 范畴 $\F(M)$ 所需的全部数据。

\StatementNumber{4}\begin{proposition}\label{st:19.4}
在拟等价意义下，$\F(M)$ 与 $c$ 及 $\phi$ 无关。
\end{proposition}
\begin{proof}
整个讨论与\secref{10d}中的讨论有些类似。第一步说明，\emph{固定 $\phi$ 时，$c$ 的选择无关紧要}。另取 $d>c$。定义 $\F(X^{\phi,d})$ 时，对于包含在较小子集 $M=X^{\phi,c}$ 中的正合 Lagrangian 子流形，可以使 Floer 数据与扰动数据只使用在 $X^{\phi,d}\setminus\operatorname{int}(X^{\phi,c})$ 的某个邻域上等于 $I_X$ 的殆复结构，以及支撑于 $\operatorname{int}(X^{\phi,c})$ 中的 Hamilton 函数，即非齐次项。最大值原理说明，非齐次伪全纯多边形留在 $X^{\phi,c}$ 内，因此 $\F(X^{\phi,c})$ 作为全 $\Ai$ 子范畴包含在 $\F(X^{\phi,d})$ 中。另一方面，$X^{\phi,d}$ 中每个正合 Lagrangian 子流形，都能由负 Liouville 流推入 $X^{\phi,c}$；在 Stein 情形中，该流就是 $\phi$ 的负梯度流。这说明包含 $\F(X^{\phi,c})\to\F(X^{\phi,d})$ 是拟等价。

第二，\emph{若 $\phi=\psi$ 在某个紧子集之外成立，则相应 Fukaya 范畴相同}。可在两个函数之间插值，再应用标准的 Moser 型论证，说明充分大的次水平集辛同构，而且该辛同构正合，并在边界附近为恒等。

第三，\emph{若如\stref{引理}{19.1}那样，从 $\phi$ 换成某个 $h(\phi)$，则 Fukaya 范畴保持不变}。可选择具有同样性质的辅助函数 $\widetilde h$，使它在 $r\gg0$ 时与 $h$ 一致，且 $r\le c$ 时 $\widetilde h(r)=ar+b$，其中 $a>0$、$b\in\R$ 为常数。从 $\phi$ 换成 $\widetilde h(\phi)$，只使次水平集 $X^{\phi,c}$ 的辛几何发生常数倍缩放，这很容易通过调整 Floer 数据与扰动数据加以处理。由于两个函数在该次水平集之外都没有临界点，相应 Fukaya 范畴相同。另一方面，从 $\widetilde h(\phi)$ 换成 $h(\phi)$ 是紧支撑的改变，因此上面的论证适用。

\SourcePage{302}
作了这些准备后，进入主要论证。给定有限型的穷竭多重次调和函数 $\phi,\psi$。由\stref{引理}{19.1}，找到比 $\psi$ 增长得快的函数 $h(\phi)$。取 $d$ 大于 $\psi$ 的全部临界值。\stref{引理}{19.2}给出 $\widetilde\phi$，它在无穷远处等于 $h(\phi)+\text{常数}$，在 $X^{\psi,d}$ 上与 $\psi$ 一致。再重复同样的论证：取 $c$ 大于 $\widetilde\phi$ 的全部临界值，且使 $X^{\psi,d}\subset\operatorname{int}(X^{\widetilde\phi,c})$。再次由\stref{引理}{19.1}，找到 $k(\psi)$，它在 $X^{\psi,d}$ 上等于 $\psi$，并比 $\widetilde\phi$ 增长得快。再应用\stref{引理}{19.2}，找到另一个函数 $\widetilde\psi$，使它在无穷远处等于 $k(\psi)+\text{常数}$，在 $X^{\widetilde\phi,c}$ 上等于 $\widetilde\phi$。

由上述讨论，$\widetilde\phi$ 相应的 Fukaya 范畴与 $\phi$ 的相同。$\widetilde\psi$ 与 $\psi$ 也如此，但在后一情形中还可作进一步观察。考虑在两个函数之间线性插值并应用 Moser 技巧所得的 $(X,\omega_{\widetilde\psi},\theta_{\widetilde\psi})$ 与 $(X,\omega_{k(\psi)},\theta_{k(\psi)})$ 之间的辛同构。该同构在 $X^{\psi,d}$ 上为恒等，因为在那里 $\widetilde\psi,\widetilde\phi,\psi,k(\psi)$ 仅相差常数。$k(\psi)$ 的负 Liouville 流将 $(X,\omega_{k(\psi)},\theta_{k(\psi)})$ 的任意紧子集推入这个次水平集。因此，应用 Moser 辛同构后，$(X,\omega_{\widetilde\psi},\theta_{\widetilde\psi})$ 的每个正合 Lagrangian 子流形都同痕于一个位于 $X^{\psi,d}$ 中的子流形。

取 $\widetilde d\gg c$。如\stref{命题}{10.5}那样，包含
\[
(X^{\widetilde\phi,c},\omega_{\widetilde\phi},\theta_{\widetilde\phi})
\hookrightarrow(X^{\widetilde\psi,\widetilde d},\omega_{\widetilde\psi},\theta_{\widetilde\psi})
\]
在 Fukaya 范畴上诱导一个在上同调上全忠实的函子。另一方面，该包含的像含有 $X^{\psi,d}$。因此，由前面的观察，该函子是拟等价。
\end{proof}

\StatementNumber{5}\begin{note}\label{st:19.5}
还有其他将 Stein 流形截取成紧子集的可能方式。例如，若 $X\subset\CC^N$ 为仿射簇，可考虑它与多球
\[
P=\bigl\{|x_1|^2+\cdots+|x_{N_1}|^2\le c_1;\ \ldots;\
|x_{N_1+\cdots+N_{r-1}+1}|^2+\cdots+|x_N|^2\le c_r\bigr\}
\]
的交 $M=X\cap P$，其中 $N_1+\cdots+N_r=N$，$c_1,\ldots,c_r$ 充分大。选一个多重次调和函数，使其负 Liouville 流相对于每个变量 $x_k$ 都指向内部，例如标准函数 $\tfrac14|x|^2$，或 Fubini--Study Kähler 势 $\log(1+|x|^2)$。相应一形式与二形式的限制使 $M$ 成为带角的正合辛流形。由\stref{命题}{10.5}，Fukaya 范畴 $\F(M)$ 与先前构造所得的相同。

另一种等价方法是放弃紧性，直接在整个 $X$ 上工作，使用无穷远处等于 $I_X$ 的殆复结构。这不完全符合本书的框架，但也许是最自然的想法，而且还能推广到无限型情形。
\end{note}
